\subsection{Sistema QUBIT}
Sistema che pu\`o esistere in due stati \(\rightarrow\) notazione di Dirac. \(\ket{A}\) ket \(\rightarrow\) stati del sistema \(\sim\) vettore.
\begin{figure}
\centering
\begin{tikzpicture}
  \node (A) at (0, 3) {};
  \node (B) at (0, 1.05) {};

  \node (C) at (0, .95) {};
  \node (D) at (0, -.95) {};

  \node (E) at (0, -1) {};
  \node (F) at (0, -3) {};

  \node (G) at (2, 3) {};
  \node (H) at (2, -3) {};

  \node (I) at (-2, 0) {\( \ket{\psi} \)};
  \node (J) at (.7, .95) {\( \ket{+} \)};
  \node (K) at (.7,-1) {\( \ket{-} \)};

  \draw[thick] (A) -- (B);
  \draw[thick] (C) -- (D);
  \draw[thick] (E) -- (F);
  \draw[thick] (G) -- (H);
\end{tikzpicture}
\caption{Nell'esperimento della doppia fenditura, lo stato \( \ket{\psi} \) \`e la combinazione degli stati \( \ket{+} \) e \( \ket{-} \).}
\label{stati_qubit}
\end{figure}
In questo caso indico gli stati con \(\ket{+}\) e \(\ket{-}\) \(\Rightarrow \ket{\pm}\) (figura \ref{stati_qubit}).

\bigskip
\noindent
Se non misuro da quale fenditura passa la particella allora ho sovrapposizione, cio\`e
\[
\ket{\psi} = a_+ \ket{+} + a_- \ket{-} \qquad \text{con } a_+, a_- \in \mathbb{C}, \mathrm{ampiezze}
\]
e \(\ket{\psi}\) \`e ancora uno stato del sistema. La probabilit\`a allora \`e definita come
\[
P_\pm = |a_\pm|^2
\]
quando sto misurando.

\bigskip
\noindent
Questa probabilit\`a \`e tale che \(P_+ + P_- = 1\), condizione di normalizzazione su \(\ket{\psi}\). A priori, \(a_\pm\) non sono determinabili.

\bigskip
\noindent
Una peculiarit\`a \`e quella di poter applicare il principio  di sovrapposizione. Quindi (senza fare misure) gli stati si sovrappongono linearmente (figura \ref{sovrapposizione_stati}).
\begin{figure}
\centering
\begin{tikzpicture}
  \node (A) at (0, 3) {};
  \node (B) at (0, 1.05) {};

  \node (C) at (0, .95) {};
  \node (D) at (0, -.95) {};

  \node (E) at (0, -1) {};
  \node (F) at (0, -3) {};

  \node (G) at (2, 3) {};
  \node (H) at (2, 1.05) {};

  \node (I) at (2, .95) {};
  \node (J) at (2, -.95) {};

  \node (K) at (2, -1) {};
  \node (L) at (2, -3) {};

  \node (M) at (-2, 0) {\( \ket{\psi} \)};
  \node (N) at (.7, .95) {\( \ket{+} \)};
  \node (O) at (.7,-1) {\( \ket{-} \)};
  \node (P) at (2.7, .95) {\( \ket{\psi_1} \)};
  \node (Q) at (2.7,-1) {\( \ket{\psi_2} \)};

  \draw[thick] (A) -- (B);
  \draw[thick] (C) -- (D);
  \draw[thick] (E) -- (F);
  \draw[thick] (G) -- (H);
  \draw[thick] (I) -- (J);
  \draw[thick] (K) -- (L);
\end{tikzpicture}
\caption{Esempio di sovrapposizione lineare degli stati.}
\label{sovrapposizione_stati}
\end{figure}
\[
\ket{\psi_1} = a_+^{(1)} \ket{+} + a_-^{(1)} \ket{-}
\]
\[
\ket{\psi_2} = a_+^{(2)} \ket{+} + a_-^{(2)} \ket{-}
\]
Sovrapposizione
\[
\ket{\psi} = c_1 \ket{\psi_1} + c_2 \ket{\psi_2}
\]
Quando misuro, sapendo che
\[
c_1 = c_2 = \frac{1}{\sqrt{2}}
\]
per semplicit\`a, e 
\[
\ket{\psi} = \frac{1}{\sqrt{2}} \left( \left( a_+^{(1)} + a_+^{(2)} \right) \ket{+} + \left( a_-^{(1)} + a_-^{(2)} \right) \ket{-} \right)
\]
ricavo le probabilit\`a
\begin{align*}
P_\pm &= \left| a_\pm^{(1)} + a_\pm^{(2)} \right|^2 = \\
&= \left| a_\pm^{(1)} \right|^2 + \left| a_\pm^{(2)} \right|^2 + \left( a_+^{(1)}\right)^* a_\pm^{(2)} + \left( a_\pm^{(2)} \right)^* a_\pm^{(1)} \neq \\
&\neq P_\pm^{(1)} + P_\pm^{(2)}
\end{align*}
Allora suppongo che ci sia un prodotto scalare: bra-ket.
\[
\ket{\psi}
\]
\[
\ket{\varphi} \rightarrow \mathrm{bra } \bra{\varphi}
\]
\[
\text{prodotto scalare} \braket{\varphi|\psi}
\]
e lo definisco con le propriet\`a
\begin{enumerate}
\item \(\braket{\varphi|\psi} = (\braket{\psi|\varphi})^*\)
\item \(\braket{\varphi|\varphi} \geq 0\)
\item se \(\ket{\psi} = a \ket{1} + b \ket{2}\) allora
\[
\braket{\varphi|\psi} = a \braket{\varphi|1} + b \braket{\varphi|2}
\]
\end{enumerate}
Cos\`{i} \(\ket{\psi} = a \ket{1} + b \ket{2}\) d\`a
\begin{align*}
\braket{\psi|w} &\stackrel{(1)}{=} \left( \braket{w|\psi} \right)^* \\
&\stackrel{(3)}{=} \left( a \braket{w|1} + b \braket{w|2} \right)^* \\
&\stackrel{(2)}{=} a^* \braket{1|w} + b^* \braket{2|w} \\
&\stackrel{(3)}{=} \left( a^* \bra{1} + b^* \bra{2} \right) \ket{w}
\end{align*}
cio\`e
\[
\bra{\psi} = a^* \bra{1} + b^* \bra{2}
\]
Stati esaustivi ed esclusivi \`e dato da
\[
\braket{+|+} = \braket{-|-}  = 1 \quad \text{esaustivo}
\]
\[
\braket{+|-} = \braket{-|+} = 0 \quad \text{esclusivo}
\]
Quindi
\[
P_\pm = \left| \braket{\pm|\psi} \right|^2 = \left| \Braket{\pm|\left( a_+ \ket{+} + a_- \ket{-} \right)} \right|^2 = |a_\pm|^2
\]
Ad esempio,
\[
\ket{\psi_1} = \frac{1}{\sqrt{2}} (\ket{+} + \ket{-})
\]
\[
\ket{\psi_2} = \frac{1}{\sqrt{2}} (\ket{+} - \ket{-})
\]
\[
P_\pm = \frac{1}{2}
\]
tuttavia non so se \(\psi_1\) e \(\psi_2\) sono uguali. Infatti,
\[
\braket{\psi_1|\psi_2} = \frac{1}{2} (\bra{+} + \bra{-})(\ket{+} - \ket{-}) = \frac{1}{2} (1 + 0 + 0 - 1) = 0
\]
cio\`e sono ortogonali, quindi sono diversi pur avendo la stessa probabilit\`a. Inoltre,
\[
\braket{\psi_1|\psi_1} = \braket{\psi_2|\psi_2} = 1
\]
Quindi se dallo stato \(\ket{\psi_1}\) misuro \(\ket{+}\) e da qui misuro la probabilit\`a associata allo stato \(\ket{\psi_2}\), classicamente otterrei 0, mentre qui ho la rigenerazione di \(\ket{\psi_2}\).
\[
\left| \braket{\psi_2|+} \right|^2 = \frac{1}{2} \left| \left( \bra{+} + \bra{-} \right) \ket{+} \right|^2 = \frac{1}{2}, \quad \text{cio\`e il 50\%}.
\]
La misura quindi modifica lo stato. Allo stesso modo, se non eseguo nessuna misura, la sovrapposizione \`e data da
\begin{align*}
\ket{\varphi} &= \frac{1}{\sqrt{2}} \left( \frac{1}{\sqrt{2}} \left( \ket{+} + \ket{-} \right) + \frac{1}{\sqrt{2}} \left( \ket{+} - \ket{-} \right) \right) = \\
&= \frac{1}{2} \left( \ket{+} + \ket{+} + \ket{-} - \ket{-} \right) = \\
&= \ket{+}
\end{align*}
cio\`e \`e sparito \(\ket{-}\) (cfr. le onde in controfase).

\bigskip
\noindent
Potrei avere
\[
\ket{\psi_A} = \frac{1}{\sqrt{2}} \left( \ket{+} + i \ket{-} \right), \qquad \ket{\psi_B} = \frac{1}{\sqrt{2}} \left( \ket{+} - i \ket{-} \right)
\]
e allora
\[
\braket{\psi_A|\psi_B} = \frac{1}{2} \left( \bra{+} - i \bra{-} \right) \left( \ket{+} - i \ket{-} \right) = \frac{1}{2} (1 - 1) = 0
\]
\[
\left| \braket{\psi_A|+} \right|^2 = \left| \frac{1}{\sqrt{2}} \left( \bra{+} - i \bra{-} \right) \ket{+} \right|^2 = \frac{1}{2}
\]
\[
\ket{\varphi} = \frac{1}{\sqrt{2}} \left( \ket{\psi_A} + \ket{\psi_B} \right) = \frac{1}{2} \left( \ket{+} + i \ket{-} + \ket{+} - i \ket{-} \right) = \ket{+}
\]
proprio come prima.

\bigskip
\noindent
Quindi se ho due stati
\[
\ket{\psi} = a_+ \ket{+} + a_- \ket{-}
\]
\[
\ket{\varphi} = b_+ \ket{+} + b_- \ket{-}
\]
allora
\[
\braket{\varphi|\psi} = \left( b_+^* \bra{+} + b_-^* \bra{-} \right) \left( a_+ \ket{+} + a_- \ket{-} \right) = a_+ b_+^* + a_- b_-^*
\]
e questo \`e leggibile in termini di sovrapposizione
\[
\ket{w} = N \left( c \ket{\varphi} + d \ket{\psi} \right)
\]
Allora \(\left| \braket{\varphi|\psi} \right|^2\) \`e una probabilit\`a che ci sia la sovrapposizione degli stati.

\bigskip
\noindent
Un altro esempio \`e
\[
\ket{\psi} = \frac{1}{\sqrt{2}} \left( \ket{+} + \ket{-} \right)
\]
\[
\ket{\varphi} = \frac{1}{\sqrt{2}} \left( \ket{+} + e^{i \vartheta} \ket{-} \right)
\]
\[
\left| \braket{\varphi|\psi} \right|^2 = \frac{1}{4} \left| 1 + e^{-i \vartheta} \right|^2
\]
Se faccio la misura
\[
\left| \braket{\varphi|+} \right|^2 = \frac{1}{2}
\]
\`e sparita la fase, quindi \`e cambiato lo stato.